% TOP_PROBLEM: {"id": 30006929, "title": "Chen-Yang Volume Conjecture for Turaev-Viro Invariants", "category_id": 7, "category": {"id": 7, "name": "topology", "display_name": "Topology", "description": "Properties preserved under continuous deformations.", "slug": "topology", "order_index": 7, "created_at": "2026-07-31T15:26:25.671Z"}, "status": "open"}
% FIRST_POSTED: 2026-09-27
% ATTEMPT_STATUS: unresolved
% !TeX program = lualatex
\documentclass[11pt]{article}
\usepackage[margin=25mm]{geometry}
\usepackage{fontspec}
\setmainfont{Latin Modern Roman}
\usepackage{amsmath,amssymb,mathtools}
\usepackage{unicode-math}
\setmathfont{Latin Modern Math}
\usepackage{xurl,hyperref}
\hypersetup{colorlinks=true,urlcolor=blue}
\setcounter{secnumdepth}{0}
\setlength{\parindent}{0pt}
\setlength{\parskip}{5pt}
\setlength{\emergencystretch}{3em}
\begin{document}
\section*{\texorpdfstring{Chen-Yang Volume Conjecture for Turaev-Viro Invariants}{Chen-Yang Volume Conjecture for Turaev-Viro Invariants}}\label{30006929}
\subsection{Definitions and mathematical statement}
Let $M$ be a finite-volume hyperbolic three-manifold in the closed, cusped, or totally geodesic boundary class specified in Chen--Yang's Conjecture 1.1. Its hyperbolic volume $\operatorname{Vol}(M)$ uses sectional curvature $-1$. For odd integers $r$, $\operatorname{TV}_r(M;q)$ is the Turaev--Viro quantum invariant: a finite sum over admissible labels of a triangulation using quantum dimensions and quantum $6j$-symbols. For cusps and boundary, use precisely the compact-core extension and normalization in the source.

The assertion is
\[
 \lim_{\substack{r\to\infty\\r\text{ odd}}}
 \frac{2\pi}{r}\log\operatorname{TV}_r(M;e^{2\pi i/r})
 =\operatorname{Vol}(M).
\]
The choice of root of unity is part of the problem. The logarithm and invariant normalization are those of the stated source; changing roots or replacing this invariant by a different quantum invariant changes the target. This is an asymptotic growth statement, not equality at finite level.
\par\smallskip\textit{Formulation and concept references:} \href{https://arxiv.org/html/1503.02547}{[S1]}.

\subsection{Short English statement}
Does the exponential growth rate of these Turaev--Viro invariants recover hyperbolic volume for every manifold in the stated class?
\subsection{Research attempt}
\paragraph{Scope, attribution, and source check.}
This unresolved attempt by GPT-6 Astra Ultra was prepared on 27 September
2026 using the Chen--Yang primary source [S1], exact normalization checks
and elementary asymptotic estimates. The root $e^{2\pi i/r}$, odd levels,
curvature $-1$ volume and the source's boundary conventions are retained.
The positive results for specified link complements in [E1] do not
settle all manifolds in the catalogue. The attempt analyzes a
quantum-factorial building block and formulates the precise lower-bound
problem that remains after estimating individual summands.

\paragraph{Polynomial factors do not affect the exponential rate.}
Fix a triangulation with $E$ edges. At level $r$ the number of edge
labelings is at most $r^E$, and admissibility only reduces this count.
If a positive normalization factor $p_r$ satisfies
$r^{-C}\leq p_r\leq r^C$, then
\[
                       \frac{2\pi}{r}\log p_r\longrightarrow0.
\]
Thus polynomial counting and normalization factors can be controlled
without knowing the exact leading coefficient of the invariant.

This observation does not authorize arbitrary renormalization.
A factor $e^{cr}$ changes the limiting rate by $2\pi c$.
Nor can one replace the root of unity by a different one:
it changes the quantum-number factors themselves and can
change their exponential behavior. All estimates below keep
the source's root fixed.

\paragraph{An elementary quantum-factorial asymptotic.}
For odd $r$ and $0\leq k\leq(r-1)/2$, put
\[
 F_r(k)=\prod_{j=1}^k2\sin(2\pi j/r),\qquad F_r(0)=1,
 \qquad
 \Lambda(\theta)=-\int_0^\theta\log|2\sin u|\,du.
\]
No factor in this range vanishes. Uniformly in $k$,
\begin{equation}\label{a431:factorial}
 \frac{2\pi}{r}\log F_r(k)
                 =-\Lambda(2\pi k/r)+O(\log r/r).
\end{equation}
To prove this, apply a Riemann-sum comparison to
$f(t)=\log(2\sin 2\pi t)$ on $(0,1/2)$.
The function increases up to $1/4$ and decreases afterwards.
On either monotone portion the difference between the
right-endpoint sum and $r$ times the integral is bounded
by endpoint oscillations. The first interval adjoining
zero is handled separately using
$f(t)=\log t+O(1)$ there; its error is bounded.
At the other end the last possible sample is at
distance $1/(2r)$ from $1/2$, so its absolute
value is $O(\log r)$. Splitting at the maximum,
with at most one crossing cell, gives
\[
 \sum_{j=1}^kf(j/r)=r\int_0^{k/r}f(t)\,dt+O(\log r).
\]
The change of variable $u=2\pi t$ proves
\eqref{a431:factorial}. The logarithmic endpoint
singularities have been controlled rather than
discarded by an unjustified uniform-continuity claim.

In the convention
$[j]_q=(q^j-q^{-j})/(q-q^{-1})$ at this root,
\[
                  [k]_q!=
             \frac{F_r(k)}{(2\sin(2\pi/r))^k}.
\]
For a product of factorials and inverse factorials
whose signed indices sum to zero, the denominator
powers cancel exactly. Formula
\eqref{a431:factorial} then gives a finite sum
of Lobachevsky terms with error $O(\log r/r)$,
provided the number of factors is fixed and
all indices lie in the displayed range.
Other index ranges or quantum-symbol conventions
require their own sign and zero-factor checks.
This building-block calculation is not yet an
asymptotic formula for a complete $6j$-symbol.

\paragraph{A Laplace principle with all hypotheses visible.}
Consider finite label sets $K_r$ in a compact set $K$,
with $|K_r|\leq r^C$, positive weights $W_r(x)$,
and a continuous action $S:K\to\mathbb R$. Suppose
\[
 \sup_{x\in K_r}
 \left|\frac{2\pi}{r}\log W_r(x)-S(x)\right|
                                  \longrightarrow0.
\]
Suppose also that for each sufficiently large odd $r$
there is $x_r\in K_r$ tending to a maximizer of $S$.
Then
\begin{equation}\label{a431:laplace}
 \lim_{\substack{r\to\infty\\r\ {\rm odd}}}
       \frac{2\pi}{r}\log\sum_{x\in K_r}W_r(x)
                                      =\max_KS.
\end{equation}
For the upper bound use the largest term times $|K_r|$;
its extra normalized logarithm is at most
$2\pi C\log r/r$. For the lower bound retain the
single term $W_r(x_r)$. These two estimates prove
the claim and show exactly where positivity is used.

A statement only along some subsequence of odd levels
would not give the target limit. Admissible labels
must approach a maximizer at every sufficiently
large odd level in the argument just given.
Uniformity near boundary labels also matters:
pointwise asymptotics at each fixed shape do not
control a maximum that moves toward a singular
boundary as $r$ grows.

\paragraph{The geometric maximization in one ideal tetrahedron.}
For an ideal hyperbolic tetrahedron with positive
dihedral angles $\alpha,\beta,\gamma$ summing to $\pi$,
the established volume formula [E2] is
\[
 V(\alpha,\beta)=
 \Lambda(\alpha)+\Lambda(\beta)+\Lambda(\pi-\alpha-\beta).
\]
Its derivatives are explicit:
$\partial_\alpha V=\log(\sin\gamma/\sin\alpha)$,
and similarly for $\beta$. Its negative Hessian is
\[
 \begin{pmatrix}
 \cot\alpha+\cot\gamma&\cot\gamma\\
 \cot\gamma&\cot\beta+\cot\gamma
 \end{pmatrix}.
\]
The first diagonal entry is
$\sin\beta/(\sin\alpha\sin\gamma)>0$, and its
determinant equals one, using
$\alpha+\beta+\gamma=\pi$.
Thus $V$ is strictly concave on the angle triangle.
The only interior critical point is
$\alpha=\beta=\gamma=\pi/3$. The volume extends
to zero at the boundary, so this is its unique
maximum, of value $v_3=3\Lambda(\pi/3)$.

This is an exact geometric check of a potential
action, not a claim that every quantum $6j$-symbol
has this particular domain or this upper bound.
Quantum-symbol asymptotics can involve different
polyhedra and different angle regimes. A proof
for a manifold must identify the actual action
and impose its gluing constraints; adding
independently maximized tetrahedral volumes
would generally ignore those constraints.

\paragraph{Cancellation is the first missing lower-bound mechanism.}
A state sum may have signed or complex summands even
when the total invariant is nonnegative in an
applicable normalization. For general summands
$z_{r,x}$ the triangle inequality gives an upper
bound from $\max|z_{r,x}|$, but no lower bound.
The two terms
\[
                    e^{rA},\qquad -e^{rA}+1
\]
both have exponential size $e^{rA}$ for $A>0$
while their sum is one. Positivity of the
final sum does not fix this cancellation.
Squaring an amplitude also does not undo
cancellation already present inside it.

A sufficient repair would be an estimate
\[
 \left|\sum_xz_{r,x}\right|
               \geq e^{-o(r)}\max_x|z_{r,x}|,
\]
together with the correct action asymptotic.
Alternatively one could find a genuinely
positive expansion to which
\eqref{a431:laplace} applies. Neither is
established here for arbitrary manifolds.
This is why knowing the largest factorial
ratio or a stationary critical shape does
not by itself prove the conjecture.

For a proposed positive expansion there are
still geometric obligations: identify its
maximal action with the complete hyperbolic
volume, verify admissibility at every large
odd level, and rule out other label regimes
with larger action. A formal critical-point
equation alone proves none of these global
comparisons.

\paragraph{Verification, gap, and outcome.}
The factorial normalization and its balanced
cancellation were checked at finite odd
levels, including indices near both
logarithmic endpoints. The tetrahedral
Hessian identity is exact; numerical values
only supplement it. No computed Turaev--Viro
values or finite range of levels is used
as proof of an asymptotic limit.

The first missing input is a uniform
asymptotic and noncancellation mechanism
for the actual invariant in the stated
manifold classes, followed by the correct
global geometric maximization.
\textbf{UNRESOLVED.} The building-block
asymptotic, conditional Laplace principle
and ideal-tetrahedron maximization are
proved. They do not establish the
Chen--Yang limit in general or justify
a change of root or normalization.

\subsection{Sources}
\begin{itemize}
\item[S1]Qingtao Chen and Tian Yang, Volume conjectures for the Reshetikhin-Turaev and the Turaev-Viro invariants, v4 (2018).
\url{https://arxiv.org/html/1503.02547}.
\textit{Formulation source.}
\item[E1]Renaud Detcherry, Efstratia Kalfagianni and Tian Yang, \emph{Turaev--Viro invariants, colored Jones polynomials and volume}.
\url{https://arxiv.org/abs/1701.07818}.
\textit{Primary source for specified link-complement results, not unrestricted volume convergence; consulted 27 September 2026.}
\item[E2]John W. Milnor, \emph{Computation of volume}, Chapter 7 of William Thurston's \emph{The Geometry and Topology of Three-Manifolds}, 1980 notes, electronic edition 1.1 (2002), especially Theorem 7.2.1.
\url{https://library.slmath.org/nonmsri/gt3m/PDF/7.pdf}.
\textit{Primary text for the ideal-tetrahedron volume formula; consulted 27 September 2026. The Hessian and maximum are computed explicitly above.}
\end{itemize}
\end{document}
